\section{Discussion}
\label{sec:discussion}

\subsection{The Tracking Limit Is Surprisingly High}

The most striking finding is that frontier 2025 models maintain above 93\% accuracy even at 64~characters.
This is a large improvement over prior benchmarks: \citet{kim2024code} reported 64.9\% for Code Llama~70B on the 7-entity \boxes task, and \citet{locations2025} found only 62\% accuracy on character-location tracking in fairy tales.
Several factors likely contribute to this gap.
First, our synthetic stories use a rigid and predictable structure (introduction $\rightarrow$ events $\rightarrow$ query) that may be easier to parse than naturalistic narratives with complex discourse structure.
Second, the structured JSON output format may prompt more careful tracking than free-form answers.
Third, the GPT-4.1 family represents a substantial capability jump over the models tested in prior work.

\subsection{Why Confusion, Not Forgetting?}

The dominance of confusion errors (93.9\% for \gptmini) has implications for how transformers represent tracked entities.
If models maintained independent per-character memory slots, we would expect errors to be primarily \emph{forgetting}---a slot's content is lost or corrupted, and the model either returns the initial value or gives up.
Instead, the model retrieves a plausible attribute value that happens to belong to the wrong character.
This is consistent with character attributes being stored in overlapping distributed representations in the transformer's residual stream, where increased entity count leads to interference between representations.

This behavioral pattern aligns with the mechanistic findings of \citet{li2025state}, who identified an Associative Algorithm in which transformers compute state updates through a parallel scan over attention heads.
In such an architecture, ``memory'' for different entities shares the same representational substrate.
As more entities compete for bandwidth, the probability of cross-entity interference increases---exactly the confusion pattern we observe.

\subsection{Why Does \gptmini Outperform \gptfull?}

Counterintuitively, the smaller \gptmini achieves higher overall accuracy (95.3\% vs.\ 93.6\% across all extended-experiment conditions).
We identify two contributing factors.

First, \gptfull has a specific weakness with compound possession-change sentences: it correctly processes the ``drop'' event but fails to register the subsequent ``pickup.''
This accounts for the majority of its errors at 32~characters and produces the anomalous dip in its degradation curve.
\gptmini does not exhibit this pattern.

Second, optimized smaller models may allocate attention more efficiently for structured tracking tasks.
The entity-tracking task has a relatively simple information structure---each character's state depends only on their own events---and may not benefit from the additional capacity that makes \gptfull superior on more complex reasoning tasks.
However, we cannot distinguish these explanations without access to model internals.

\subsection{Gradual Degradation Implies Continuous Capacity}

We observe no sharp ``cliff'' in accuracy as character count increases.
Instead, accuracy degrades smoothly from 100\% at 16~characters to $\sim$93\% at 64~characters.
This gradual pattern is more consistent with a continuous resource (representational bandwidth shared across entities) than a discrete slot model (where we would expect accuracy to remain high until slots are exhausted, then drop sharply).

\subsection{Limitations}

Our study has several limitations that bound the generalizability of our findings.

\para{Synthetic stories only.}
Our narratives follow a rigid three-phase structure with explicit state changes and unambiguous attribute assignments.
Real narratives involve implicit state changes, coreference, dialogue, unreliable narration, and complex discourse structure.
Our accuracy numbers likely represent an upper bound on what models can achieve in naturalistic settings.

\para{Two models from one provider.}
We test only GPT-4.1 and GPT-4.1-mini.
Results may not generalize to other model families (Claude, Gemini, Llama, Qwen).

\para{Two attribute types.}
We track only locations and possessions---concrete, easily verifiable attributes.
Real narrative understanding requires tracking beliefs, intentions, emotions, and relationships, which may be substantially harder.

\para{Structured output format.}
The JSON batch-query format may prompt more careful tracking than open-ended questions.
An interactive question-by-question protocol might yield different results.

\para{Small sample sizes at high $N$.}
The extended experiment uses only 3~stories per character count.
While confidence intervals account for sampling variability, the true variability across story structures may be larger.

\para{Temperature 0 does not eliminate stochasticity.}
OpenAI models are not fully deterministic even at temperature~0 due to non-deterministic GPU operations, so our results may vary slightly across runs.
